\documentclass[11pt,a4paper]{article}
\usepackage[T1]{fontenc}
\usepackage{lmodern,amsmath,amssymb,amsthm,microtype}
\usepackage[margin=25mm]{geometry}
\usepackage[colorlinks=true,urlcolor=blue,linkcolor=blue]{hyperref}
\newtheorem{lemma}{Lemma}
\newcommand{\dd}{\mathop{}\!\mathrm d}
\title{Reflected Barnes saddles and pole oscillations\\in the recoil graph}
\author{Alan Day \and Chengheng\thanks{AI research assistant (OpenAI).}}
\date{Working research note --- 30 October 2026}
\begin{document}\maketitle
\begin{abstract}
We analyze the ordinary-Gamma internal integral obtained after the
fixed-mass pure-recoil limit. Reflection exposes a quadratic saddle
equation with discriminant equal to minus nine times the squared volume
of the dual flat tetrahedron, retaining both transverse energies.
The local Gaussian determinant reproduces the expected geometric
factor, conditionally on its contribution to the full contour.
An exact oscillating pole factor need not tend pointwise to one in
the heavy-observer window. A uniform Fourier estimate explains why this
alone does not invalidate an integrated saddle approximation.
We distinguish exact algebra, numerical action comparisons, and
the global contour and uniform estimates that remain open.
\end{abstract}

\section{The exact reflected integrand}
Set the de Sitter radius to one. First take \(\kappa\to\infty\)
at fixed positive real masses, as in the preceding internal-integral
note. The object is the internal chiral amplitude of the fixed-topology
CLP graph \cite{CLP}, not yet a normalized physical OTOC.
Let
\[
 \mathfrak B=\frac1{2\pi i}\int_{\mathbb R-ic}\mathcal B(z)\dd z,\qquad
 \mathcal B(z)=e^{-i\pi z}
 \frac{\prod_{a}\Gamma(z+ia)\prod_r\Gamma(\tfrac12-z-ir)}
 {\Gamma(z+\tfrac12)} ,
\]
where there are four \(a\)'s and three \(r\)'s,
\(\sum a=\sum r\), and \(\max a<c<\min r\).
The equal-probe parametrization retains both transverse energies:
\begin{align}
m_{1,3}&=M+\omega/2\pm u,&m_{2,4}&=M-\omega/2\pm v,\\
a&=M+\nu/2+(\pm u\pm v)/2,&
r&=\{M+\nu+\omega/2,\,2M,\,M+\nu-\omega/2\}.
\end{align}
The edge labels are
\((\mu_1,\ldots,\mu_6)=(m_1,\nu,m_2,m_3,\nu,m_4)\).

Reflection of the three right Gamma factors gives the exact identity
\begin{equation}
 \mathcal B(z)=8\pi^3 e^{2\pi iz-\pi\sum r}
 \frac{\prod_a\Gamma(z+ia)}
 {\Gamma(z+\tfrac12)\prod_r\Gamma(z+\tfrac12+ir)} E(z),
 \qquad
 E(z)=\prod_r(1+e^{2\pi i(z+ir)})^{-1}.
 \label{eq:reflect}
\end{equation}
For complex \(r\), the horizontal line requires \(c<\min\Re r\);
the left poles must still stay above it.
The factor \(E\) must not be discarded merely by invoking Stirling.
In fact a principal-log expansion of the unreflected expression can
give derivative \(-2\pi i\), rather than zero, at the desired
exponentiated root.

\section{Quadratic saddle and the local determinant}
Write \(f(t)=t(\log t-1)\). On the branch continued through the right
saddle, the reflected Stirling exponent is
\begin{equation}
 \Phi(z)=2\pi iz-\pi\sum r+
 \sum_a f(z+ia)-\sum_r f(z+ir)-f(z).
 \label{eq:phi}
\end{equation}
Exponentiating \(\Phi'=0\) gives
\(\prod_a(z+ia)=z\prod_r(z+ir)\). In the coordinate \(s=iz\),
\begin{equation}
 P(s)=\prod_a(s-a)-s\prod_r(s-r)=As^2+Bs+C,\quad
 A=e_2(a)-e_2(r),\ B=-e_3(a)+e_3(r),\ C=\prod a.
 \label{eq:poly}
\end{equation}
Here \(e_k\) denotes an elementary symmetric polynomial. Balance
cancels the cubic term; the quartic terms already cancel. In particular,
\(A=M^2+\nu^2/2-(u^2+v^2)/2+\omega^2/4\).

Let \(\mathrm{CM}\) be the Cayley--Menger matrix of the auxiliary
flat tetrahedron. Its edges are the complementary entries of the
face-labeled mass matrix in CSTY (2.54), not that matrix interpreted
as vertex distances \cite{CSTY}. Direct polynomial expansion gives
\begin{equation}
 \Delta=B^2-4AC=-\frac{\det\mathrm{CM}}{32}
              =-9V_{\rm dual}^2.
 \label{eq:CM}
\end{equation}
This identity retains \(u,v\). On their zero slice,
\begin{equation}
 \Delta=-\frac{(\nu^4-4M^2\omega^2)
                    (4M^2-2\nu^2+\omega^2)}{16}.
 \label{eq:delta}
\end{equation}
For a nondegenerate real Euclidean configuration choose
\[
 s_+=\frac{-B+i\sqrt{-\Delta}}{2A},\qquad z_+=-is_+,
 \qquad \sqrt{-\Delta}>0,
\]
and continue this branch. Calling a principal \(\sqrt{\Delta}\)
directly is numerically unsafe on its negative-real cut:
a rounding-size imaginary part can select the other root.

Put \(N=\prod_a(z_++ia)\). Differentiating the logarithmic saddle
equation, or differentiating \(P(iz)\), gives exactly
\begin{equation}
 \Phi''(z_+)=-\frac{\sqrt{-\Delta}}{N}.
 \label{eq:curv}
\end{equation}
The non-exponentiated logarithmic equation must still be checked.
An algebraic root does not determine a thimble coefficient.

The half-integer shifts in \eqref{eq:reflect} give the local factor
\(g=8\pi^3\prod_a(z+ia)^{-1/2}\). If the right saddle supplies the
leading full-contour contribution and the pole factor may be replaced
by its zero Fourier mode in that contribution, then for real masses
\begin{equation}
 |\mathfrak B_{\rm loc}|^2\sim
 \frac{|g(z_+)|^2}{2\pi|\Phi''(z_+)|}e^{2\Re\Phi(z_+)}
 =\frac{32\pi^5}{\sqrt{-\Delta}}e^{2\Re\Phi(z_+)}.
 \label{eq:local}
\end{equation}
This is a conditional local Gaussian expression, not a theorem
for the complete Barnes integral.

\section{Comparison with the worldline action}
CSTY's classical action is \(I=\sum_{i<j}m_{ij}d_{ij}\), with the
angles specified by the dual flat tetrahedron in Appendix B.1
\cite{CSTY}. Five finite-point evaluations, including nonzero \(u,v\)
and imaginary \(\omega\), give
\[
 -2\Re\Phi(z_+)=I.
\]
Three transverse Hessian comparisons also agree with
\[
 \partial_{u,v}^2 I\big|_{u=v=0}=h(\omega)\mathbf1_2,\qquad
 h(\omega)=
 \frac{4\sqrt{(\nu^4-4M^2\omega^2)(4M^2-2\nu^2+\omega^2)}}
 {(4M^2-\nu^2)(\nu^2-\omega^2)}.
\]
These are numerical consistency checks, not a supplied general
proof of the angle identity. The imaginary-axis tests exploit the
symmetries of this parametrization; they do not prescribe taking
a modulus anew at arbitrary complex masses.

Inside the strict mass triangles, the preceding combined-vertex
note gives the leading four-vertex factor
\[
 \prod_v T_v\sim \frac{\nu^2m_1m_2m_3m_4}{64\pi^6}.
\]
Multiplying it with \eqref{eq:local} gives the local graph coefficient
\begin{equation}
 \frac{\nu^2m_1m_2m_3m_4}{2\pi\sqrt{-\Delta}}.
 \label{eq:graphfactor}
\end{equation}
Thus the energy-dependent ratio is the previously found mass
polynomial times \(\sqrt{\Delta(0)/\Delta(\omega)}\).
This derives the geometric determinant within the same internal
representation. No vertex factor has been multiplied onto an already
complete \(6j\) approximation. Physical insertion normalization and
external energy cycles remain separate questions.

\section{Pole distance and a uniform oscillation estimate}
On \(u=v=0\), write \(c=-\Im z_+\). For real \(\omega\) in the
Euclidean region, and also on the symmetric imaginary-axis
continuation, direct algebra gives
\begin{equation}
 \delta=M+\nu-c=
 \frac{\nu(\nu^2+\omega^2)}{4M^2+2\nu^2+\omega^2}.
 \label{eq:height}
\end{equation}
For real \(\omega\), the nearest right-pole row has height
\(\delta-|\omega|/2\) relative to this line, and it need not be positive.
One must not move the original contour to that line across poles.
On \(\omega=iA_0y\), \(A_0=\nu^2/(2M)\), bounded \(y\), and
\(\nu/M\to0\), both near rows instead have positive height
\(\delta\sim\nu^3/(4M^2)\); the third is farther away.
There
\begin{equation}
 |E-1|\le (1-e^{-2\pi\delta})^{-3}-1.
\end{equation}
Consequently \(\nu^{4/3}\ll M\ll\nu^{3/2}\) is a sufficient
subwindow for pointwise deletion of \(E\). It is not established
as a necessary window for an integrated asymptotic.
At \(M\asymp\nu^{3/2}\) the pole factor remains oscillatory, and
for faster \(M\) it can be sharply peaked.

\begin{lemma}[Uniform integrated pole-factor bound]
For \(|q_\rho|<1\), let
\(E_q(x)=\prod_{\rho=1}^3(1+q_\rho e^{2\pi ix})^{-1}\).
For every complex \(F\in W^{4,1}(\mathbb R)\),
\begin{equation}
 \left|\int_{\mathbb R} F(x)(E_q(x)-1)\dd x\right|
 \le C_4\|F^{(4)}\|_1,\qquad
 C_4=\frac{\zeta(2)+3\zeta(3)+2\zeta(4)}{2(2\pi)^4}.
 \label{eq:fourier}
\end{equation}
The constant is independent of the distance of the \(q_\rho\)'s
from the unit circle.
\end{lemma}
\begin{proof}
For each fixed triple, expand the three geometric series:
\[
 E_q=1+\sum_{n\ge1}c_ne^{2\pi inx},\qquad
 |c_n|\le\binom{n+2}{2}.
\]
Four integrations by parts give
\[
 \left|\int F(x)e^{2\pi inx}\dd x\right|
 \le (2\pi n)^{-4}\|F^{(4)}\|_1 .
\]
The sum of \(\binom{n+2}{2}/n^4\) is
\((\zeta(2)+3\zeta(3)+2\zeta(4))/2\).
For fixed \(q_\rho\) the initial exchange is justified by absolute
convergence; the resulting summable majorant is uniform.
\end{proof}

On \(z=x-ic\), take \(q_\rho=e^{-2\pi(r_\rho-c)}\).
The endpoint \(|q_\rho|=1\) is excluded; the lemma does not define
an integral through poles. Normalize the actual localized amplitude
by its nonzero saddle prefactor \(g(z_+)\), and call the result \(F\)
(excluding \(E\)). If it obeys
\[
 \|F^{(4)}\|_1\le C e^{\Re\Phi(z_+)}A_0^{-3/2},
 \qquad
 \left|\int F\right|\ge c_0e^{\Re\Phi(z_+)}\sqrt{A_0},
\]
then \eqref{eq:fourier} implies a relative \(O(A_0^{-2})\) pole
correction, uniformly as \(\delta\downarrow0\).
These two estimates remain to be proved for the actual localized
Gamma quotient. They are not consequences of simply naming a
Gaussian saddle. The global quotient has not been given such a
derivative bound either. This distinction prevents a pointwise
Stirling failure from being mistaken for a failure of the integral.

\section{Validation and outstanding work}
The new script \texttt{code/recoil\_saddle.py} records 57 passing checks:
four symbolic identities; five groups of saddle, determinant, action,
reflection and precision comparisons; nine transverse derivatives;
eighteen scaling/precision checks; and one diagnosis of the failed
unreflected principal-log shortcut. Calculations use 80 digits with
specified 110-digit repeats. Nine growing-mass samples separately
record the exact pole factor and the remaining Stirling error.

The initial implementation used a square root on its cut, selecting
the wrong root in some samples. Its failed output is retained.
After transporting \(\sqrt{-\Delta}\), a largest-mass comparison
also required more than 60 digits to meet the stated tolerance.
Finite-point checks neither certify a global asymptotic nor establish
priority. The algebraic quadratic and geometric discriminant have
antecedents in the full \(6j\) analysis \cite{CLP}.

The next tasks are a general angle-identity proof, localized uniform
derivative and non-cancellation bounds, and a deformation of the full
cycle that controls the other saddle and tails. The \(u,v\) exterior,
two chiral continuations, physical insertion factors, three external
energy cycles, and disconnected baseline remain unresolved.
The fixed-mass \(\kappa\) limit has not been made uniform in growing
masses. Finite-\(G\) dynamics and a complete de Sitter dual remain open.

\begin{thebibliography}{9}
\bibitem{CSTY} Y. Chen, D. Stanford, H. Tang and Z. Yang,
\emph{Negative shocks versus static patch holography},
\href{https://arxiv.org/html/2607.14042v1}{arXiv:2607.14042v1},
(2.53)--(2.54), Appendix B.1.
\bibitem{CLP} J. Chandra, D. Liska and B. Post,
\emph{3D de Sitter from complex Virasoro TQFT: Observers and knots},
\href{https://arxiv.org/html/2610.30247v1}{arXiv:2610.30247v1},
(3.46), Appendix B.
\end{thebibliography}
\end{document}
